\documentclass[10pt,a4paper]{article}
\usepackage[T1]{fontenc}\usepackage[utf8]{inputenc}
\usepackage{lmodern,amsmath,amssymb,booktabs,array,microtype}
\pdfmapfile{+lm.map}\pdfmapfile{+cm.map}\pdfmapfile{+cmextra.map}\pdfmapfile{+symbols.map}
\pdfinfoomitdate=1\pdftrailerid{}
\usepackage[a4paper,margin=22mm]{geometry}
\usepackage{xcolor,hyperref,fancyhdr}
\definecolor{ink}{HTML}{183A4A}
\hypersetup{colorlinks=true,linkcolor=ink,urlcolor=ink,pdftitle={CST-DUST: Gewichtsrobustheit des vollstaendigen Blocks v0.1}}
\pagestyle{fancy}\fancyhf{}\fancyhead[L]{\small X6 | CST-DUST Gewichtsrobustheit}
\fancyhead[R]{\small v0.1 | 12.11.2026}\fancyfoot[C]{\thepage}
\setlength{\headheight}{14pt}\setlength{\parindent}{0pt}\setlength{\parskip}{5pt}
\setlength{\emergencystretch}{2em}\renewcommand{\arraystretch}{1.15}
\newcommand{\A}{\mathcal A}\newcommand{\D}{\mathcal D}\newcommand{\one}{\mathbf1}
\newcommand{\W}{\mathcal W}\newcommand{\dd}{\mathrm d}
\begin{document}
{\color{ink}\LARGE\bfseries Gewichtsrobustheit des vollständigen Blocks}\par
{\large Globales Minimum über die bisher zugelassene Historiengewichtsklasse}\par
\textbf{Ergebnis: Eine vollständige Auslöschung ist in dieser Klasse ausgeschlossen.}

Für dieselben 33 AA-Historien, dieselben sechs Ein- und Ausgangstestzustände,
dieselbe $V$-Norm und dieselbe feste Entkopplung wie im vorigen Schritt sei
\[
 \Delta_6(w)=R_{F_6}\bigl(K_1(w)-K_0(w)\bigr)|_{E_6},\qquad
 m_6=\min_{w\in\W}\|\Delta_6(w)\|_{V\to F_6'}.
\]
$\W$ enthält \emph{alle} nichtnegativen, normierten, randinvarianten Gewichte.
Die globale Einschließung lautet
\[
 \boxed{0{,}0022939794998\ \leq\ m_6\ \leq\ 0{,}0022939795002.}
\]
Die linke Schranke gilt für jedes zulässige Gewicht. Die rechte Schranke
wird durch ein explizites rationales Gewicht bezeugt. Beide Schranken
enthalten Quadraturrest, Rundung und Zertifikatsfehler; sie sind keine
statistischen Konfidenzintervalle.

Auch der gesamte $4\times2$-Atom$\to$Positiv-Block besitzt eine positive
Untergrenze. Seine \emph{separate} Minimierung ergibt
\[
 \boxed{0{,}0003502439596\ \leq\ m_{0\to+}\ \leq\ 0{,}0003502439599.}
\]
Die beiden beinahe minimierenden Gewichte sind verschieden. Es wird kein
gemeinsames Minimum beider Ziele behauptet.

\begin{center}\begin{tabular}{lrr}\toprule
Norm des Differenzblocks & $6\times6$ & Atom$\to$Positiv\\\midrule
Gleichmäßige gelabelte Gewichte & 0,032602301976 & 0,012323446639\\
Bisherige Einzelkanal-Auslöschung & 0,008446681427 & 0,001702265450\\
Globales Minimum, näherungsweise & 0,002293979500 & 0,000350243960\\\bottomrule
\end{tabular}\end{center}
Die ersten beiden Zeilen reproduzieren den vorigen Befund; die letzte
Zeile fasst die obigen Intervalle zusammen. Die Werte sind normierte
Amplituden, keine Wahrscheinlichkeiten. Das Maß beeinflusst ihre Größe
deutlich, kann aber den hier geprüften Gesamtbefund nicht auf null drücken.

\textbf{Der zusammenhängende Historienanteil bleibt gesondert.}
Für ihn existiert keine positive uniforme Untergrenze: Bei Konzentration
auf einen einzelnen Randorbit verschwindet er exakt. Bei strikt positiven
Gewichten auf allen Historien kann er beliebig klein werden. Der robuste
Differenzbefund darf deshalb nicht als gewichtsrobuster Nachweis eines
zusammenhängenden Historiensignals umgedeutet werden.

\textbf{Geltungsbereich.} Zertifiziert ist ein endlicher Test-/Antidualblock.
Ein physikalisch ausgezeichnetes CST-Historienmaß, eine volle
$V$-Operatornorm, physikalische Evolution und ein Bounce folgen daraus nicht.

\newpage
\section{Unveränderte Klasse, Zustände und Normierung}
Die zuvor deklarierte Klasse lautet ausdrücklich: nichtnegative, normierte
Gewichte, invariant unter unabhängigen Vertauschungen der zwei Eingangs-
und Ausgangslabels, unabhängig von $\lambda$ und Integrationsvariablen.
Der eingefrorene Vertrag steht in \texttt{inputs/weight\_class\_contract.json}.

Seien $o=1,\ldots,10$ die Orbits von $S_2(L)\times S_2(R)$ mit Größen $n_o$.
Die gesamte Klasse wird exakt parametrisiert durch
\[
 \W=\{w\in\mathbb R^{10}:w_o\geq0,\ \sum_o w_o=1\},\qquad
 c_C=w_o/n_o\quad(C\in o).
\]
Es gibt \textbf{keine} zusätzliche Untergrenze für $w_o$, keine Bindung an
gleichmäßige Gewichte, keine Entropieauflage und keine aus dem Ergebnis
abgeleitete Verklebungsbedingung. Die Orbitreduktion folgt aus der schon
festgelegten Randsymmetrie. Sie fügt keine neue physikalische Maßannahme hinzu.

Die sechs Zustände sind $\Psi_{ab}=\sqrt2\,d_a h_b$ mit konstantem Uhrpaket,
\[
 d_0=2\one_{00},\quad d_+=2(\one_{10}+\one_{01}+\one_{11})/\sqrt3,\quad
 d_c=2(\one_{10}+\one_{01}-2\one_{11})/\sqrt6,
\]
\[
 h_0=\pi^{-1/2}e^{-|q|^2/2},\qquad h_R=(|q|^2-1)h_0.
\]
$E_6=F_6$ ist ihr Spann; die Reihenfolge ist $(0,0),(0,R),(+,0),(+,R),(c,0),(c,R)$.
Für den Teilblock sind $E_2=\operatorname{span}(\Psi_{0,0},\Psi_{0,R})$
und $F_4=\operatorname{span}(\Psi_{+,0},\Psi_{+,R},\Psi_{c,0},\Psi_{c,R})$.

Die Norm ist unverändert $V=L^2_{\rm inv}((1+|q|^2)^4\dd\mu_A)$ mit dem
zuvor festgelegten Randmaß. Die Dichte-/Uhrfaktoren sind orthonormal; je
Feldpaar gilt exakt
\[
 G=\begin{pmatrix}65&196\\196&913\end{pmatrix},\quad \det G=20929,\qquad
 T=\begin{pmatrix}1/\sqrt{65}&-196/\sqrt{65\cdot20929}\\
 0&\sqrt{65/20929}\end{pmatrix},\quad T^*GT=I.
\]
Alle folgenden Matrizen werden mit dieser \emph{vollen} Gram-Matrix normiert.
Ein Basiswechsel ändert ihre Operatornorm nicht, wenn die Grams mitwechseln.

Die Intervention bleibt $M_{C,\lambda}=\overline M+\lambda(M_C-\overline M)$,
$\lambda=0,1$, mit $\overline M=(1/33)\sum_C M_C$. Diese Referenz wird bei
Gewichtsänderungen nicht neu gemittelt. Gravitation, Staub und Randpakete
bleiben an beiden Endpunkten gleich. Insbesondere werden die Amplituden
nicht durch die komplexe Gravitationssumme dividiert.

\section{Lineare Blockdarstellung und exakte Kompression}
Seien $\A_o$ und $\A_{\rm ref}$ die analytischen $2\times2$-Feldantworten,
$d_o\in\mathbb C^3$ die eingeschlossenen Dichteintegrale und
$p_o=e^{iF_4(o)}$. Definiere den $6\times2$-Block $B_o$ durch
\[
 (B_o)_{2a:2a+2,:}=p_o(d_o)_a\,T^*(\A_o-\A_{\rm ref})T,\qquad a=0,1,2.
\]
Diese Beiträge enthalten alle inneren Dichtemasken und die vollständigen
Gravitationsphasen. Sie werden genau einmal berechnet; $w_o$ bleibt außen.

\newpage
\section{Ein globales Zertifikat statt eines lokalen Optimierergebnisses}
Mit $e=\int_0^1e^{-i\rho/2}\dd\rho$ lautet die feste Eingangsdichtezeile
\[
 r=(1,(2e+e^2)/\sqrt3,(2e-2e^2)/\sqrt6),\quad
 R=(r_0 I_2\;r_+ I_2\;r_c I_2),\quad J=R^*/\|r\|_2.
\]
Dann ist $J^*J=I_2$ und der vollständige normierte Block erfüllt
\[
 \widehat\Delta_6(w)=\Bigl(\sum_o w_oB_o\Bigr)R
 =C(w)J^*,\qquad C(w)=\sum_o w_oC_o,\quad C_o=\|r\|_2B_o.
\]
Also gilt \emph{exakt} $\|\widehat\Delta_6(w)\|_2=\|C(w)\|_2$.
Die Kompression entfernt keine empfindliche Eingangsrichtung. Vier
Dichte-Nullrichtungen bleiben für jedes Gewicht bestehen. Für den
Atom$\to$Positiv-Block werden stattdessen nur die unteren vier Zeilen von
$B_o$ verwendet; dort kommt der Faktor $\|r\|_2$ nicht hinzu.

Das Problem $\min_{w\in\W}\|\sum_o w_oC_o\|_2$ ist konvex [2]. Zur
Zertifizierung genügt hier sogar ein einziger fester Test- und Eingangsvektor.
Seien $u\ne0$, $v\ne0$ die gespeicherten \emph{rational komplexen}
Kandidaten. Für jeden Orbit wird mit Ballarithmetik eingeschlossen:
\[
 q_o=\frac{\operatorname{Re}(u^*C_ov)}{\|u\|_2\|v\|_2}.
\]
Für \textbf{jedes} zulässige $w$ folgt direkt aus Cauchy--Schwarz
\[
 \left\|\sum_o w_oC_o\right\|_2
 \geq\frac{\operatorname{Re}\bigl(u^*(\sum_o w_oC_o)v\bigr)}{\|u\|_2\|v\|_2}
 =\sum_o w_oq_o\geq\min_o q_o.
\]
Die Rechnung beweist $\min_o q_o>0{,}0022939794998680$ für den Gesamtblock.
Die publizierte untere Endstelle $0{,}0022939794998$ liegt sicher darunter.
Für den Teilblock gilt analog $\min_o q_o>0{,}0003502439596714$.
Der volle Zeuge hat Eingang $Jv/\|v\|_2$ und Test $u/\|u\|_2$. Somit
existiert sogar \emph{ein festes Paar von $V$-Einheitszuständen}, dessen
Realteil für alle Gewichte positiv bleibt. Seine Koeffizienten stehen im Zertifikat.

Für einen rational zulässigen Kandidaten $w^{\rm rat}$ wird die obere
Schranke unabhängig berechnet. Ist $H=C(w^{\rm rat})^*C(w^{\rm rat})$
und $H=\left(\begin{smallmatrix}a&z\\\bar z&d\end{smallmatrix}\right)$, so gilt
\[
 \sigma_{\max}(C)=\left[\frac{a+d+\sqrt{(a-d)^2+4|z|^2}}2\right]^{1/2}.
\]
Alle Ausdrücke werden nach außen eingeschlossen. Es ergibt sich
$\sigma_{\max}(C)<0{,}0022939795000484$ und für den Teilblock
$\sigma_{\max}<0{,}0003502439598197$.

Ein numerischer Optimierer liefert lediglich die Kandidaten. Weder sein
Abbruchkriterium noch seine Zielfunktionswerte gehen als Annahme in den
Beweis ein. Die zehn separierenden Ungleichungen und ein zulässiges
rationales Gewicht reichen aus, um das \emph{globale} Minimum einzuschließen.

\newpage
\section{Explizite Gewichte und strikt positiver Vollsupport}
Die folgenden Zahlen sind exakte endliche Dezimalzahlen, keine gerundete
Definition. Alle nicht aufgeführten Orbitmassen sind null. Das gelabelte
Gewicht ist jeweils $c_C=w_o/n_o$.
\begin{center}\begin{tabular}{rrrr}\toprule
Orbitcode & $n_o$ & $w^{\rm rat}_6$ & $w^{\rm rat}_{0\to+}$\\\midrule
8476 & 4 & 0 & 0,0126227577805495\\
8592 & 4 & 0,2084830921717939 & 0,2174457610437640\\
8604 & 4 & 0,3092377970167804 & 0\\
8988 & 4 & 0,2704525417920602 & 0,2443946647275104\\
24860 & 4 & 0,2118265690193655 & 0,5255368164481761\\\bottomrule
\end{tabular}\end{center}
Beide Spalten summieren sich exakt zu eins. Sie bezeugen die oberen
Schranken; sie werden nicht als exakt bestimmte Minimierer ausgegeben.

Die Klasse strikt positiver Gewichte ist im Simplex relativ offen. Für sie ist zunächst ein
\emph{Infimum}, nicht zwingend ein angenommenes Minimum, zu untersuchen.
Da die Norm stetig ist und jede Randgewichtung durch
\[
 w_o^{(\varepsilon)}=(1-\varepsilon)w_o^{\rm rat}
                       +\varepsilon\frac{n_o}{33},\qquad 0<\varepsilon<1,
\]
mit strikt positiven Gewichten approximiert werden kann, gilt allgemein
\[
 \inf_{w_o>0,\,\sum w_o=1}\|\Delta(w)\|
 =\min_{w_o\geq0,\,\sum w_o=1}\|\Delta(w)\|.
\]
Für die konkrete Wahl $\varepsilon=10^{-12}$ sind schon alle 33 gelabelten
Gewichte positiv. Die verifizierten oberen Schranken sind
\[
 \|\Delta_6(w^{(\varepsilon)})\|<0{,}0022939795000504,\qquad
 \|\Delta_{0\to+}(w^{(\varepsilon)})\|<0{,}0003502439598199.
\]
Auch diese Werte liegen innerhalb der auf Seite 1 angegebenen Intervalle.
$\varepsilon$ ist ein \emph{Zeugenparameter}; er ist keine neue zulässige
Mindestgewichtung und wird der Klasse nicht als Bedingung auferlegt.

\section{Was das Minimum über das Historienmaß aussagt}
Die gleichmäßige gelabelte Gewichtung liegt deutlich über dem globalen
Minimum. Die Stärke des endlichen Befundes ist deshalb maßabhängig.
Seine Nichtnull-Eigenschaft ist über die \emph{gesamte erklärte Klasse}
robust. Insbesondere können auch die zuvor gefundenen positiven
Einzelkanal-Auslöschungsgewichte den gesamten Block nicht auslöschen.

Dies unterscheidet drei Aussagen: Ein einzelner Eintrag kann null sein;
der vollständige Block kann zugleich ungleich null sein; und der neue
Beweis liefert eine positive Untergrenze für die Blocknorm über alle
zulässigen Gewichte. Er liefert keine positive Antwort für \emph{jeden}
Eingangszustand. Die vier festen Nullrichtungen des Gesamtblocks bleiben bestehen.

Für ein anderes Historienuniversum, andere Randzustände, eine andere
$V$-Norm oder eine mit $w$ veränderte Referenzwirkung müsste die Rechnung
neu durchgeführt werden. Ein physikalisches Zähl- oder Verklebungsgesetz
für die Gewichte wurde hier weder benötigt noch hergeleitet.

\newpage
\section{Zusammenhängender Anteil und mittlere Skalarantwort}
Die aus dem vorigen Schritt übernommene Diagnose verwendet
\[
 p_o(w)=w_oe^{iF_4(o)},\quad Z(w)=\sum_o p_o(w),\qquad
 \widetilde\A(w)=\frac{\sum_o p_o(w)\A_o}{Z(w)}.
\]
Nur wo $Z(w)\ne0$ ist diese Zerlegung definiert. Dann gilt, vor identischer
$V$-Normierung beider Seiten,
\begin{align*}
 \Delta(w)&=R_{\rm conn}(w)+R_{\rm mean}(w),\\
 R_{\rm conn}(w)&=\sum_o p_o(w)(d_or)\otimes(\A_o-\widetilde\A(w)),\\
 R_{\rm mean}(w)&=\left(\sum_o p_o(w)d_or\right)
                        \otimes(\widetilde\A(w)-\A_{\rm ref}).
\end{align*}
Die Gewichte $p_o/Z$ sind komplex und keine Wahrscheinlichkeiten.
Der primäre Differenzblock bleibt auch ohne diese Diagnose definiert;
Gewichte werden in der Minimierung \emph{nicht} wegen $Z=0$ ausgeschlossen.

\begin{center}\begin{tabular}{lrrr}\toprule
Gewicht, gesamter $6\times6$-Bereich & $\|\Delta\|$ & $\|R_{\rm conn}\|$ & $\|R_{\rm mean}\|$\\\midrule
Gleichmäßig gelabelt & 0,032602301976 & 0,007251598185 & 0,030294483160\\
$w^{\rm rat}_6$ & 0,002293979500 & 0,002251686356 & 0,001730771896\\
$w^{\rm rat}_{0\to+}$ & 0,003218220579 & 0,002447005160 & 0,003959407301\\\bottomrule
\end{tabular}\end{center}
Jede Tabellenzahl hat einschließlich Rundung absoluten Fehler $<10^{-12}$.
Die Matrizen addieren sich exakt; ihre Normen sind nicht additiv.

\textbf{Exakter Gegenbefund für den zusammenhängenden Anteil.}
Für ein Orbitvertex $w=e_o$ ist $Z=e^{iF_4(o)}\ne0$ und
$\widetilde\A=\A_o$. Daher gilt exakt
\[
 R_{\rm conn}(e_o)=0,\qquad R_{\rm mean}(e_o)=\Delta(e_o)\ne0.
\]
Das Gewicht ist randinvariant und gehört zur ursprünglichen Klasse.
Die skalare Antwort ist innerhalb des Orbits konstant, auch wenn mehrere
Labels enthalten sind. Für $w^{(\varepsilon)}=(1-\varepsilon)e_o+
\varepsilon(n_1,\ldots,n_{10})/33$ bleiben alle Gewichte positiv und
$Z$ ist nahe $\varepsilon=0$ von null getrennt. Stetigkeit liefert
$R_{\rm conn}(w^{(\varepsilon)})\to0$. Also ist auch über strikt positiven
Gewichten das Infimum des zusammenhängenden Anteils null.

\section{Fehlerbilanz, Reproduktion und Quellen}
Die zehn Orbitintegrale werden aus dem vorigen Paket unverändert übernommen
und mit 40 Gauß--Legendre-Knoten pro Uhrdimension und 192 Bit vollständig
neu reproduziert. Die positiven Dichten sind analytisch integriert;
Bernstein-Ellipsenreste schließen die Uhrquadratur ein. Die Feldintegration
verwendet den analytischen Abel-Grenzwert mit kontrollierten Gaußphasen.
Arb/Acb propagiert Rest- und Rundungsintervalle bis zu jeder Dualungleichung
und zur oberen Singularzahlschranke. Das gesamte Gewichts-Simplex wird
durch zehn Ungleichungen erfasst; ein Gewichts-Rasterfehler existiert nicht.

\texttt{python -B verify.py} prüft die Hashes, rechnet in einer temporären
Kopie alle Orbitintegrale und Ergebnisdateien neu und verlangt identische
numerische Ergebnisse. Neun Kontrollgruppen enthalten fünf gezielte
Negativkontrollen. Kandidatensuche: \texttt{discover.py}; sie ist kein Teil
des Zertifikatsvertrauens. Python-Versionen stehen in \texttt{environment.json}.

\small
[1] CST-DUST, Basisunabhängiger Kopplungsbefund v0.1, Main-Commit\\
\texttt{558339f0c40a3309d108674d411f548d2061dd57},
\href{https://github.com/thomasmgfeller/x6-spectral-astrolabium/pull/287}{PR 287}.
Gewichtsklasse: Mischungsqualifikation v0.1, eingefrorener Vertrag im Paket.

[2] S. Boyd, L. Vandenberghe, \emph{Convex Optimization}, Cambridge University
Press, 2004, Beispiel 3.11 (Konvexität der Spektralnorm),
\href{https://web.stanford.edu/~boyd/cvxbook/}{offizielle Buchseite}.
Der hier verwendete Separationsbeweis ist oben vollständig angegeben.

Neue Arbeit: Apache-2.0. Kein physikalisch ausgezeichnetes Maß, keine externe
Begutachtung, kein neuer 1089-Historienlauf. Kompakte Reduktion bleibt HOLD;
physikalische CST-Evolution bleibt NOT\_CONSTRUCTED.
\end{document}
